\section{Related Work}
\label{sec:related_work}

\para{Literature-grounded idea generation.} Early LLM ideation studies showed that models can generate plausible research ideas, but the best systems increasingly add literature grounding. \textsc{SciMON} retrieves inspirations from prior papers and iteratively revises ideas to reduce overlap with literature \citep{wang2024scimon}. \textsc{ResearchAgent} uses scientific literature and reviewing agents to iteratively generate research problems, methods, and experiments \citep{baek2024researchagent}. \textsc{SciPIP} improves retrieval with semantic, citation, and full-text information, then combines retrieved content with LLM internal knowledge \citep{wang2024scipip}. \textsc{Chain of Ideas} organizes papers into chains that model how research directions develop over time \citep{li2024chain}. Unlike these systems, we do not propose a full ideation agent. We isolate prompt-level grounding operations under matched topics, model, and evaluation criteria.

\para{Evaluation of scientific ideas.} Evaluation is the central challenge for automated ideation. \citet{si2024llms} ran a controlled human study with more than 100 NLP researchers and found that LLM ideas can be judged more novel than expert ideas while being weaker on feasibility and diversity. \rinobench turns expert review data into a benchmark for novelty judgment and shows that LLM rationales can resemble expert rationales while numeric novelty scores still diverge from human labels \citep{schopf2026rinobench}. The Idea Novelty Checker frames novelty assessment as a literature-grounded retrieval and reranking problem \citep{shahid2025literature}. We use an independent LLM judge because the present experiment evaluates newly generated ideas, but we report automatic overlap and paired tests to avoid relying only on one scalar score.

\para{Idea-generation benchmarks.} \ideabench and \aibidea provide structured paper and reference contexts for research idea generation \citep{guo2024ideabench,qiu2025aiideabench}. These benchmarks are useful because they separate a research context from a target paper idea. We use \aibidea contexts as input topics, but we do not expose hidden target-paper methods, motivations, or summaries to the generator. Our goal is not to recover the target paper; it is to estimate how prompt grounding changes the quality and literature proximity of generated ideas.

\para{Execution as grounding.} A stronger form of grounding is to run experiments. \textsc{MLAgentBench} evaluates agents on machine-learning experimentation tasks \citep{huang2023mlagentbench}; \textsc{DiscoveryBench} formalizes data-driven discovery tasks \citep{majumder2024discoverybench}; and \textsc{ScienceAgentBench} tests scientific agents on validated coding tasks from scientific workflows \citep{chen2024scienceagentbench}. \textsc{The AI Scientist} combines idea generation, code execution, paper writing, and simulated review \citep{lu2024aiscientist}. Our \smallexperiment condition is much weaker: it uses a real but shallow lexical audit over titles. The negative result for this condition supports the broader lesson that grounding needs meaningful evidence, not just an empirical-looking step.
